\section{Furstenberg set problem: upper bound}\label{sec3}

In this section, we prove Theorem \ref{upper bound}. In other words, we will construct a $(s,t;n,k)$-set $E=\bigcup_{V\in\cV}Y(V)$, such that
\[ \#E\lesssim_{s,t,n,k} p^{\bF(s,t;n,k)}.\]
  

\subsection{Two special cases} To familiarize ourselves with $\bF(s,t;n,k)$, we discuss two special cases: $s=k$; $k=1$.

When $s=k$, we can morally view $Y(V)$ as $V$. Therefore, a $(k,t;n,k)$-set is of the form
\[ E=\bigcup_{V\in\cV}V. \]
Estimating the size of $E$ is known as the problem of ``union of $k$ planes". It has been resolved in finite fields by R. Oberlin \cite{oberlin2016unions}. The Euclidean version was resolved by the author in \cite{gan2023hausdorff}.

When $k=1$, this is the Furstenberg set problem for lines. To construct a $(s,t;n,1)$-set 
\[ E=\bigcup_{L\in\cL}Y(L) \]
with size as small as possible, we hope that the lines in $\cL$ are compressed in a low-dimensional space. 
Write $t=2m+\beta$, where $0\le m\le n-2$ is an integer, and $\beta\in (0,2]$. Since $t\le 2(m+1)$, we can make $\cL$ contained in $\Fp^{m+2}$. Therefore, our goal becomes to construct a small $(s,t;m+2,1)$-set. The trick is to take the Cartesian product of a $(s,\beta;2,1)$-set with $\Fp^m$. Let $E'\subset \Fp^2$ be a $(s,\beta;2,1)$-set and let $E'$ have the form
\[ E'=\bigcup_{L'\in \cL'}Y(L'). \]
We will construct a $(s,t;m+2,1)$-set $E$ based on $E'$. Let $\pi_{\Fp^2}:\Fp^{m+2}\rightarrow \Fp^2$ be the projection onto the first two coordinates. Define $\cL$ to be the set of lines $L$ in $\Fp^{m+2}$ such that $\pi_{\Fp^2}(L)$ is a line in $\cL'$. Then, $\#\cL\sim \# G(1,\Fp^{m+1})\cdot \#\cL'\sim p^{2m+\beta}=p^t$. If $L\in \cL$, then $L'=\pi_{\Fp^2}(L)\in \cL'$. We define $Y(L):=L\cap \pi_{\Fp^2}^{-1}(Y(L'))$. We see that $\#Y(L)=\#Y(L')\sim p^s$, so 
\[ E=\bigcup_{L\in\cL}Y(L) \]
is a $(s,t;m+2,1)$-set.
Also, since $E\subset E'\times \Fp^m$, we have
\[ \#E\le\#E'\cdot p^m\lesssim p^{m+\min\{s+\beta,\frac32s+\frac12\beta,s+1\}}  \]
By checking Definition \ref{deffurindex}, one exactly sees that
    \[ m+\min\{s+\beta,\frac32s+\frac12\beta,s+1\}=\bF(s,t;m+2,1)=\bF(s,t;n,1). \]

\subsection{Proof of Theorem \ref{upper bound}}\hfill

For admissible $(s,t;n,k)$, we will construct a $(s,t;n,k)$-set 
\[ E=\bigcup_{V\in\cV}Y(V) \]
such that
\[ \#E\gtrsim p^{\bF(s,t;n,k)}. \]

\fbox{Case (a): \texorpdfstring{$s=0$}{}}

If $t\le k(n-k)$, we choose $\cV$ to be $p^t$ $k$-planes pass through the origin $\bzz^n$, and $Y(V)=\{\bzz^n\}$. Then $\#E\le 1=p^0$. Here, $\bzz^n$ means $(\underbrace{0,\dots,0}_{n \textup{~times~}})$.

If $t\ge k(n-k)$, we choose $p^{t-k(n-k)}$ points $X$ in $\Fp^{n-k}$. For each such point $x\in X$, we choose $\sim p^{k(n-k)}$ $k$-planes that pass through $x$, and require the $k$-planes to be transverse to $\Fp^{n-k}$. We obtain $\sim p^t$ $k$-planes which is $\cV$. If $V$ passes through $x\in X$, we set $Y(V)=\{x\}$. We see that $\#E=\#X\le p^{t-k(n-k)}$.

\bigskip

\fbox{Case (b): $s>0, t\in [0,(k-d-1)(n-k)]$}

We write $s=d+\si$, where $d\in \{0,\dots,k-1\}, \si\in(0,1]$.
The construction is similar to Case (a). We choose $\sim p^t$ many $k$-planes that contain $\Fp^{d+1}$. This is doable since 
\[ \#\{ V\in A(k,\Fp^n): V\supset L \}=\#G(k-d-1,\Fp^{n-d-1})\sim p^{(k-d-1)(n-k)}. \]
Let $\cV$ be these $k$-planes. Fix $Y\subset \Fp^{d+1}$ with $\#Y=p^s$. For each $V\in \cV$, let $Y(V)=Y$. We see that $\#E=\#Y=p^s$.





\bigskip

\fbox{Case (c): $s\in (0,k], t\in ((k-d-1)(n-k),(k+1)(n-k)]$}

We write 
\begin{equation}\label{stformula}
    \begin{cases}
        s=d+\si\\
        t=(k-d-1)(n-k)+(d+2)m+\tau,
    \end{cases}
\end{equation} 
    where $d\in\{0,\dots,k-1\}$, $\si\in (0,1]$,  $m\in\{0,\dots,n-k-1\}$, $\tau\in [0,2]$. We remark that here we allow $\tau=0$, although we assumed $\tau>$ in the canonical expression. Allowing $\tau=0$ will produce the example in Case (d).

We write $\Fp^n=\Fp^{n-k+d+1}\times \F_p^{k-d-1}$. Note that $\gtrsim 1$ fraction of $k$-planes intersect $\Fp^{n-k+d+1}$ at a $(d+1)$-plane. Heuristically speaking, to make $\#E$ as small as possible, we would like $E$ to be a subset of $\F_p^{n-k+d+1}\times \bzz^{k-d-1}$. For simplicity, we write $\F_p^{n-k+d+1}\times \bzz^{k-d-1}$ as $\Fp^{n-k+d+1}$, thinking of it as a subspace of $\Fp^n$. This will not cause any ambiguity.

Note that for each $W\in A(d+1,\F_p^{n-k+d+1})$, the cardinality of 
\[\{V\in A(k,\Fp^n): V\supset W,  V \textup{~is~transverse~to~} \F_p^{n-k+d+1} \}\] 
is $\sim \#G(k-d-1,\Fp^{n-k})\sim p^{(k-d-1)(n-k)}$. We just need to find $\cW\subset A(d+1,\F_p^{n-k+d+1})$ with $\#\cW\sim p^{t-(k-d-1)(n-k)}=p^{(d+2)m+\tau}$, and for each $W\in\cW$, to find $Y(W)\subset W$ with $\#Y(W)\sim p^{s}$, such that
\begin{equation}\label{constructW}
    \#\bigcup_{W\in\cW} Y(W)\lesssim p^{\bF(s,t;n,k)}.
\end{equation}

If this is done, then we let $\cV$ be those $k$-planes $V$ that contain some $W\in\cW$, and $Y(V)=Y(W)$ if $V\supset W$. We see that $\#\cV\sim \#\{W\}\cdot p^{(k-d-1)(n-k)}\sim p^t$, and hence $E=\bigcup_{V\in\cV} Y(V)$ is a $(s,t;n,k)$-set. Also, we have the upper bound
\[ \#E=\#\bigcup_{W\in\cW} Y(W)\lesssim p^{\bF(s,t;n,k)}. \]

It suffices to find a $(s,(d+2)m+\tau;n-k+d+1,d+1)$-set $E_1=\bigcup_{W\in\cW} Y(W)$ such that \eqref{constructW} holds. 
Note that $\#A(d+1,\Fp^{m+d+2})\sim p^{(d+2)(m+1)}$ and $(d+2)m+\tau\le (d+2)(m+1)$. We will construct the $(d+1)$-planes $\cW$ so that they lie in $\Fp^{m+d+2}$. 
It suffices to find a $(s,(d+2)m+\tau;m+d+2,d+1)$-set $E_1=\bigcup_{W\in\cW} Y(W)$ such that \eqref{constructW} holds.

We write $\Fp^{m+d+2}=\Fp^{d+2}\times\Fp^{m}=\Fp^2\times \Fp^d\times \Fp^m$. We first choose a $(\si,\tau;2,1)$-set $E_2=\bigcup_{L\in\cL}Y(L)$ in $\Fp^2$, such that
\[\#E_2\lesssim p^{\min\{\si+\tau,\frac32\si+\frac12\tau,\si+1\}}. \]
The existence of $E_2$ was shown in Proposition \ref{prop2d}.

Next, we define $\cU:=\{L\times \Fp^d: L\in\cL\}$ which is a subset of $A(d+1,\Fp^{d+2})$. And for each $U=L\times \Fp^d\in\cU$, define $Y(U)=Y(L)\times \Fp^d$. We see that $\#\cU= \#\cL\sim p^\tau$, $\#Y(U)=\#Y(L)\cdot p^d\sim p^{\si+d}=p^s$. We also see that $E_3:=\bigcup_{U\in \cU}Y(U)$ is contained in $E_2\times \Fp^d$, and hence \[\#E_3\lesssim p^{d+\min\{\si+\tau,\frac32\si+\frac12\tau,\si+1\}}.\]

Finally, we define $\cW$. For each $U\in \cU$, consider the $(m+d+1)$-plane $U\times \Fp^m$.
Let $\pi_U:U\times \Fp^m\rightarrow U$ be the projection (naturally defined). We let $\cW_U$ be the set of $(d+1)$-planes $W$ contained in $U\times \Fp^m$ such that $\pi_U(W)=U$. Note that $\gtrsim 1$ fraction of the $(d+1)$-planes in  $U\times \Fp^m$ lie in $\cW_U$, so $\#\cW_U\sim \#A(d+1,\Fp^{m+d+1})\sim p^{(d+2)m}$.

We define $\cW=\bigcup_{U\in \cU}\cW_U$ (actually this is a disjoint union). For $W\in \cW_U$, define $Y(W)=W\cap \pi_U^{-1}(Y(U))$. We see that $\#\cW\sim \#\cU\cdot p^{(d+2)m}\sim p^{(d+2)m+\tau}$, and $\#Y(W)=\#Y(U)\sim p^s$. Hence, $E_1=\bigcup_{W\in\cW}Y(W)$ is a $(s,(d+2)m+\tau;m+d+2,d+1)$-set. Since $E_1\subset E_3\times \Fp^m$, we have
\[ \#E_1\le \#E_3\cdot p^m\lesssim p^{d+m+\min\{\si+\tau,\frac32\si+\frac12\tau,\si+1\}}. \]
The construction is done.

\bigskip
\fbox{Case (d)}

We still use \eqref{stformula}, but the range of $\tau$ is $\tau\in(2,d+2]$. Note that the worst case is when $\tau=d+2$.
What we need to do is for
\begin{equation}
    \begin{cases}
        s=d+\si\\
        t=(k-d-1)(n-k)+(d+2)(m+1),
    \end{cases}
\end{equation} 
construct a $(s,t;n,k)$-set $E$ such that
\[ \#E\lesssim p^{s+m+1}. \]
The trick is to still use the construction in Case (c), with $\tau$ replaced by $0$ and $m$ replaced by $m+1$. Therefore, we can construct a $(s,t;n,k)$-set $E$ such that
\[\#E\lesssim p^{d+m+1+\min\{\si+0,\frac32\si+\frac120,\si+1\}=p^{d+m+1+\si}=p^{s+m+1}}. \]
\
